\chapter{\abstractname}
Fast compilation is crucial in contexts where code generation occurs at runtime, such as query compilers in modern database systems or runtimes for modern programming languages. However, the quality of initially generated code is often markedly suboptimal, prompting many systems to adopt multi-tiered compilation strategies.

Recently, TPDE has shown that compilation speed of LLVM unoptimized back-end can be improved substantially while keeping a similar runtime. 
However, because TPDE further reduces compilation time at the unoptimized level, it also widens the compile-time gap to the next optimization level, \texttt{-O1}, to two orders of magnitude.

In this work, we present a novel register allocation approach built on top of TPDE that enables improved code generation while preserving a fast compilation time. By combining two analysis passes with a repair mechanism applied during code generation, our approach optimizes spill placement and resolves instruction constraints.

On SPECint 2017 we achieve a 24\% runtime improvement over the LLVM -O0 back-end, while maintaining a 4x faster compile time on optimized LLVM IR. Although, this also results in a 5x compile time slowdown compared to the original TPDE. Our results demonstrate the potential for a faster, lightly optimizing compiler back end; however, several design and implementation issues limit the performance of our approach.
